\documentclass[../main.tex]{subfiles}
\begin{document}
\begin{center}
    \begin{overpic}[height=1.3\pagewidth]{../../../../Image/Book?/chaptitle/intro0.png}
    \put(50, 20){               % x = 50 (mặc định trong quyển này), y = 80
        \begin{tcolorbox}[
            colframe=oldpaper,  % Màu viền
            boxrule=1pt,      % Độ dày viền
            arc=0mm,         % Góc vuông (không bo tròn)
            enhanced jigsaw,
            opacityback = 0,
            opacityframe = 0,
            left=5pt,        % Khoảng cách giữa nội dung và viền trái
            right=5pt,       % Khoảng cách giữa nội dung và viền phải
            top=5pt,         % Khoảng cách giữa nội dung và viền trên
            bottom=5pt,      % Khoảng cách giữa nội dung và viền dưới
            width = 13cm     % (mặc định)
        ]
        {\color{oldpaper}
        Tháng Tư, năm 2021.
        
        \vspace{10pt}
    
        Hikawa Kuon đây. Cùng với Hikawa Kaigou-chan và Hikawa Suzuki - hai chị em đến từ một thế giới song song, và cùng với Game-san - những người bạn đồng hành của tôi kể từ khi bị lạc ở The Void.

        \vspace{5pt}
        
        Đúng, thật sự là ``Hikawa'', không phải ``Haragen'' đâu.
    
        \vspace{5pt}
    
        Chúng tôi từng bị cuốn vào một cơn lốc dữ dội và trôi dạt đến một thế giới kỳ lạ, nơi mà Game đặt mật danh là ``hololive ERROR''. Ở đó, chúng tôi mất trí nhớ, bị ép mang những cái tên giả, sống giữa những vòng lặp vô tận của động đất, ngày học lặp đi lặp lại, và cả những người bạn học - mang dáng dấp của những nàng hề ở thế giới gốc của tôi - bị chính thế lực vô hình nào đó thao túng. Từng ngày, từng đêm, chúng tôi cố gắng tìm cách thoát ra, trong khi chính mình cũng dần không phân biệt nổi đâu là thật, đâu là giả.
    
        \vspace{5pt}
    
        Và rồi, ở trung tâm của tất cả là cô gái ấy: Misora Shino. Người vừa là nạn nhân, vừa là nguồn gốc của mọi vòng lặp.

        \vspace{5pt}

        Cả ba chúng tôi cùng với Game, Shino-san, và Shionomiya Sakura-san, đã cùng nhau cứu rỗi cô ấy. Kết quả là, thị trấn Aogami được tái sinh, và tất cả các bạn học từ các vòng lặ trước đó đều được đưa trở lại tới thế giói thực.
        
        \vspace{5pt}

        \dots Ít nhất, là ở trên bề mặt.

        \vspace{5pt}
    
        Vấn đề nảy sinh từ đây: không có dấu hiệu nào cho thấy lối thoát đã mở ra. Không có lỗ xoáy, không có khe nứt không gian, không có cánh cổng nào có thể trở về thế giới cũ.
        
        Theo như Game tính toán, \textbf{vẫn còn những linh hồn trẻ em - nạn nhân từ các nghi thức hiến tế xa xưa} đang lẩn khuất quanh thị trấn này. Có thể chúng chính là phần ``chưa được khôi phục'' của nơi đây.
    
        \vspace{5pt}
    
        Vậy nên, có vẻ như chúng tôi sẽ phải\dots\ sống tiếp như những học sinh trung học thực thụ. \dots Ừ thì, ``bình thường'' không hẳn là từ đúng, vì xung quanh tôi toàn là con gái, và chuyện siêu nhiên thì lúc nào cũng lấp ló ở góc lớp. Vừa học, vừa tìm cách giúp những linh hồn đó được giải thoát, đồng thời tìm đường rời khỏi nơi này.
    
        \vspace{5pt}
    
        Dù sao thì, tôi vẫn tin rằng ở đâu đó, phía sau màn sương của Aogami, vẫn còn điều gì đó chờ chúng tôi khám phá. Và có lẽ, chìa khóa để trở về chính là \emph{khôi phục} nốt phần còn thiếu của thế giới này.
        }\end{tcolorbox}
    }
    \end{overpic}
\end{center}
\end{document}